思考分析力学被发明出来的原因
\begin{itemize}
    \item 系统有很多约束，解决多约束问题
    \item 力学和光学统一起来？
\end{itemize}
这几个公式不知道对不对：
$\sqrt{E-u\left(x\right)}=\frac{1}{f\left(\omega\right)v_{ph}\left(\omega,x\right)}$
$$\begin{cases}\frac{\mathrm{d}}{\mathrm{d}\omega}\left(\omega f\left(\omega\right)\right)=0 \\ \frac{\omega f\left(\omega\right)}{2}\frac{\mathrm{d}E\left(\omega\right)}{\mathrm{d}\omega}=\sqrt{\frac{m}{2}} \end{cases}$$
$$\sqrt{\frac{m}{2\left(E\left(\omega\right)-U\right)}}=\frac{\mathrm{d}}{\mathrm{d}\omega}\left(\omega\cdot f\left(\omega\right)\sqrt{\left(E\left(\omega\right)-U\right)}\right)$$
\chapter{薛定谔方程}
\section{波函数的统计解释和几率守恒方程}
薛定谔方程：\[
i\hbar \frac{\partial}{\partial t}\psi(\mathbf{r},t)
=
\left[
-\frac{\hbar^2}{2m}\nabla^2 + V(\mathbf{r},t)
\right]\psi(\mathbf{r},t)
\]
自己推导一下：$$\frac{\partial}{\partial t}\psi ^*\left( \vec{r} \right) \psi \left( \vec{r} \right) =\frac{i\hbar}{2m}\nabla \cdot \left[ \psi ^*\left( \vec{r} \right) \nabla \psi \left( \vec{r} \right) -\psi \left( \vec{r} \right) \nabla \psi ^*\left( r \right) \right] $$
有这个可以得到类似的密度守恒方程：$$\frac{\partial}{\partial t}\rho \left( \vec{r},t \right) +\nabla \cdot \vec{j}=0$$

定态的相加不一定为定态。